%----------------------------------------------------------------------------------------
% Tailored CV — PhD/Postdoc, In Vivo Cancer Therapeutics
% LMU University Hospital, Prof. Lars Lindner, Sarcoma / Hyperthermia / Drug Delivery
%----------------------------------------------------------------------------------------

\documentclass[10pt,a4paper,sans]{moderncv}

\usepackage[utf8]{inputenc}
\usepackage[T1]{fontenc}

\moderncvstyle{classic}
\moderncvcolor{blue}

\usepackage[scale=0.78]{geometry}
\usepackage{ragged2e}

%----------------------------------------------------------------------------------------

\firstname{Kundai Farai}
\familyname{Sachikonye}

\address{Auf der Lichtung 19}{82178 Puchheim, Germany}
\mobile{(+49) 1626386344}
\email{kundai.sachikonye@bitspark.com}
\homepage{github.com/fullscreen-triangle}{github.com/fullscreen-triangle}

%----------------------------------------------------------------------------------------

\begin{document}

\makecvtitle

%----------------------------------------------------------------------------------------
%	EDUCATION
%----------------------------------------------------------------------------------------

\section{Education}

\cventry{2019--2021}{PhD candidate, Computational Lipidomics}{Technische Universität München}{}{}{
  Mass spectrometry-based lipid systems characterisation, multi-omics statistical
  modelling, federated learning for distributed clinical analysis.
  Departed programme to pursue independent computational research.
}
\cventry{2018--2019}{MSc Computational Biology}{Universität Tübingen}{}{}{}
\cventry{2014--2017}{MSc Pharmaceutical Biotechnology}{HAW Hamburg}{}{}{
  Drug delivery, formulation science, pharmacokinetics, biopharmaceutical process
  development.
}
\cventry{2011--2014}{BSc Biochemistry and Cell Biology}{Jacobs University Bremen}{}{}{
  Cell culture, molecular biology, biochemical laboratory techniques.
}

%----------------------------------------------------------------------------------------
%	RESEARCH OUTPUT
%----------------------------------------------------------------------------------------

\section{Research Output}

\cvitem{2025}{
  \textbf{On the Thermodynamic Consequences of Categorical Completion on Biological
  Systems.}
  MHC binding affinity correlates with parent protein categorical richness
  ($r=0.73$, $p<10^{-12}$, $n=2{,}847$ IEDB peptides). AUC 0.81 vs NetMHCpan 0.68.
  Predicts sarcoma (collagen isoforms, $R \approx 10^{5.5}$) and pancreatic
  cancer proteins as primary immune targets; explains immunologically cold
  tumour phenotypes from structural first principles. \textit{Manuscript.}
}

\cvitem{2025}{
  \textbf{On Disease Epistemology in Blind Receiver.}
  Loop-holonomy self-consistency as unique viable diagnostic criterion;
  AUC$=1.00$; 25/25 validation experiments. Sparse $\ell_1$ LP for minimum-side-effect
  therapeutic target selection. \textit{Manuscript. AIMe Registry / TUM Weihenstephan.}
}

\cvitem{2025}{
  \textbf{Syndrome: Categorical Resolution of Disease Trajectories.}
  Disease vector $\mathbf{D}\in[0,1]^8$ across eight cellular oscillator classes
  (including membrane potential $D_M$ and calcium $D_{Ca}$, directly responsive
  to hyperthermia); $\mathcal{O}(N\log N)$ trajectory to target state;
  sparse LP for minimum-intervention combination therapy selection.
  \textit{Manuscript.}
}

%----------------------------------------------------------------------------------------
%	RESEARCH SOFTWARE
%----------------------------------------------------------------------------------------

\section{Research Software}

\cvitem{lavoisier}{
  \textbf{Lipidomics and Mass Spectrometry Pipeline.}
  Complete pipeline for lipid species characterisation: MS1/MS2 ion annotation,
  spectral library matching, feature extraction. Dual-pipeline architecture:
  numerical + CNN visual pipeline (PyTorch); 98.9\,\% accuracy against NIST.
  Ray distributed computing; $>$100\,GB mzML support. Applicable to liposomal
  formulation characterisation and phospholipid phase transition analysis.
  Virtual instrument: \href{https://lavoisier.vercel.app/experiment}{lavoisier.vercel.app/experiment}.
  \href{https://github.com/fullscreen-triangle/lavoisier}{github.com/fullscreen-triangle/lavoisier}
}

\cvitem{borgia}{
  \textbf{First-Principles Molecular Spectroscopy and Cheminformatics.}
  Derives atomic structure, molecular spectroscopic properties, and lipid
  phase transition temperatures from bounded phase space geometry alone ---
  no empirical parameters. The Triple Equivalence Theorem (oscillatory dynamics
  $=$ categorical state enumeration $=$ partition operations) enables prediction
  of DPPC/DSPC vibrational spectra and phase transition temperatures from
  molecular geometry. Validated: nine elements against NIST electron configurations;
  six molecules (H$_2$ through C$_6$H$_6$) with self-consistency errors $<$1.63\%.
  Live: \href{https://honjo-masamune.vercel.app/tools}{honjo-masamune.vercel.app/tools}.
  \href{https://github.com/fullscreen-triangle/borgia}{github.com/fullscreen-triangle/borgia}
}

\cvitem{syndrome}{
  \textbf{Tumour Microenvironment State, Immunopeptidology, and Therapy Optimisation.}
  Four live tools at \href{https://syndrome-seven.vercel.app}{syndrome-seven.vercel.app}
  (3-channel DFT spectral embedding, 36D cosine similarity $\rho$, $O(cL\log L)$):
  neoantigen identifier (dual threshold $\Delta\rho_\text{self}>0.15$ and
  $\rho_\text{MHC}>0.72$) for identifying which peptides released by hyperthermia-mediated
  lysis become immunogenic; antibiotic/cytotoxic resistance predictor
  (drug--enzyme spectral complementarity $\rho$, validated against TEM-1/Ampicillin,
  GyrA/Ciprofloxacin, AAC(3)-IV/Gentamicin).
  Disease vector $\mathbf{D}\in[0,1]^8$; membrane potential ($D_M$) and calcium
  ($D_{Ca}$) oscillators model bilayer thermal stress and immunogenic cell death
  signalling. Sparse $\ell_1$ LP selects minimum-intervention combination of
  hyperthermia dose, liposomal drug concentration, and immunotherapy agent.
  \href{https://github.com/fullscreen-triangle/syndrome}{github.com/fullscreen-triangle/syndrome}
}

\cvitem{gospel}{
  \textbf{Tumour Genomics and Neoantigen Analysis.}
  Processes million-scale VCF files at $10^6$ variants/second; RNA-seq integration;
  Bayesian network inference, Gaussian mixture models. Validated against gnomAD v3.1.2,
  UK Biobank. Applicable to somatic variant profiling in sarcoma and pancreatic cancer.
  Live: \href{https://vivid-symbolism.vercel.app/search}{vivid-symbolism.vercel.app/search}.
  \href{https://github.com/fullscreen-triangle/gospel}{github.com/fullscreen-triangle/gospel}
}

\cvitem{shakespear}{
  \textbf{Protein Analysis Instrument.}
  Folding, trajectory, catalysis, and dynamics via Kuramoto oscillator synchronization
  in S-entropy space. Applicable to oncoprotein structural analysis and characterisation
  of hyperthermia-released tumour antigens.
  Live: \href{https://shakespear-nine.vercel.app/instrument}{shakespear-nine.vercel.app/instrument}.
  \href{https://github.com/fullscreen-triangle/shakespear}{github.com/fullscreen-triangle/shakespear}
}

\cvitem{crown-prince}{
  \textbf{Geometric Pharmacology and Therapeutic Effect Trajectories.}
  Drug-target binding as $K_d = \exp(-\Delta M \cdot \ln b)$ from partition geometry;
  39/39 NIST validations; zero free parameters; 50\,MB vs DrugBank 12\,GB.
  Therapeutic Effect Trajectory paper derives PD/PK/ADME from first principles.
  Cell Spectral Hologram: tumour cell input $\to$ sparse $\eta^*$ therapeutic
  perturbation vector identifying minimum-cytotoxicity drug target combinations.
  Live: \href{https://crown-prince-four-courners.vercel.app/tools/spectrum}{crown-prince-four-courners.vercel.app/tools/spectrum}.
}

%----------------------------------------------------------------------------------------
%	RESEARCH EXPERIENCE
%----------------------------------------------------------------------------------------

\section{Research Experience}

\cventry{2021--present}{Independent Computational Biology Developer}{Bitspark GmbH}{}{}{
  Multi-omics analysis frameworks, first-principles molecular spectroscopy,
  tumour microenvironment modelling, and scientific data tools. All systems publicly
  deployed and validated.
}

\cventry{2019--2021}{Computational Lipidomics}{TUM / Bayerisches Forschungsinstitut}{Munich}{}{
  Mass spectrometry pipelines for lipid species characterisation (peak detection,
  ion annotation, LIPID MAPS database integration), multi-omics statistical modelling,
  federated learning for distributed clinical analysis. Python, Java.
  Instrumentation: TIMS-PASEF, QTOF, ion trap.
}

\cventry{2019}{Therapeutic Lipid Production Optimisation}{Fraunhofer Institut IGB}{Stuttgart}{}{
  Real-time process monitoring and dynamic programming / linear programming optimisation
  for industrial photobioreactor production of \textbf{therapeutic lipids}. Python, C.
}

\cventry{2017--2018}{Glycomics and Proteomics}{Universitätsklinikum Hamburg-Eppendorf}{}{}{
  Automated characterisation pipeline for LC-MS/MS and MALDI-TOF data; protein
  identification and quantification from biological samples.
}

\cventry{2013}{Cancer Biomarker Discovery}{University Medical Center Groningen}{}{}{
  UPLC-MS/MS shotgun proteomics for cancer biomarker discovery from cell plasma
  and serum; protocol optimisation and statistical analysis.
}

%----------------------------------------------------------------------------------------
%	TECHNICAL SKILLS
%----------------------------------------------------------------------------------------

\section{Technical Skills}

\cvitem{Mass Spectrometry}{
  LC-MS/MS, MALDI-TOF / MALDI MSI, UPLC-MS/MS, TIMS-PASEF, QTOF;
  lipid species annotation (LIPID MAPS, LipidBlast), phospholipid characterisation,
  spectral library matching (NIST), metabolite identification
}

\cvitem{Drug Delivery / Formulation}{
  Thermosensitive liposomal systems, phospholipid phase transitions,
  liposome characterisation by mass spectrometry; therapeutic lipid production
  and process optimisation; pharmacokinetics and drug distribution modelling
}

\cvitem{Cancer Biology}{
  Tumour microenvironment modelling, neoantigen prediction, somatic variant
  profiling, immune evasion mechanisms, CRISPR screen analysis (MAGeCK),
  pathway analysis (KEGG, Reactome, GSEA)
}

\cvitem{ML / Statistics}{
  PyTorch, scikit-learn, XGBoost; Bayesian networks, Gaussian mixture models,
  variational Bayes; survival analysis; multi-omics integration (MOFA, UMAP)
}

\cvitem{Languages}{Python, R, Rust, TypeScript, Java, C, SQL, Bash}

\cvitem{Infrastructure}{
  Ray, Dask, Docker, FastAPI, AWS, PostgreSQL, memory-mapped I/O for large datasets
}

%----------------------------------------------------------------------------------------
%	LANGUAGES
%----------------------------------------------------------------------------------------

\section{Languages}

\cvitemwithcomment{English}{Native / Fluent}{}
\cvitemwithcomment{German}{Conversational}{Living and working in Germany since 2011}

%----------------------------------------------------------------------------------------

\renewcommand{\listitemsymbol}{-~}

\end{document}
